\documentclass[11pt]{article}
\usepackage{icmlko}
\begin{document}

\title{여섯 자 중 넷이 침묵한다 --- 그리고 셋째 짝은 관문이 될 수 없다}
\author{월드모델 실험실}
\date{노트 613 · 2026년 8월 5일}
\maketitle

\begin{abstract}
노트 $611$ 이 다섯 자 중 \textbf{셋만} 표본 SD 를 냈다. 나머지 둘(날짜 통제
판 $\cdot$ 무리 안 판)과 셋째 짝(CN 만화 $352$ 행)을 같은 자로 채웠다. 이번
세션의 채택(노트 $586$)에 대해 여섯 자가 하는 말은

\begin{center}
\begin{tabular}{lrrrl}
\hline
자 & 행 & 차 & 표본 SD & $t$ \\
\hline
판 & $3{,}369$ & $+0.0015$ & $0.0021$ & $0.7$ \\
날짜 통제 판 & $3{,}369$ & $+0.0034$ & $0.0026$ & $1.3$ \\
\textbf{무리 안 판} & $3{,}369$ & $+0.0024$ & $\mathbf{0.0114}$ & $\mathbf{0.2}$ \\
\textbf{KR 만화} & $1{,}716$ & $+0.0390$ & $0.0123$ & $\mathbf{3.0}$ \\
\textbf{비게임 앱} & $1{,}600$ & $+0.2187$ & $0.0122$ & $\mathbf{17.9}$ \\
CN 만화 & $352$ & $+0.0105$ & $0.0109$ & $1.0$ \\
\hline
\end{tabular}
\end{center}

\textbf{판 계열 셋이 전부 $0$ 을 물고 CN 도 문다.} 집 밖 두 짝만 $0$ 밖이다
--- \textbf{이 채택은 그 둘 위에 서 있다.}
\end{abstract}

\begin{figure}[t]
\centering
\includegraphics[width=\textwidth]{figs/f1.pdf}
\caption{\textbf{가} --- 여섯 중 둘만 $0$ 밖이다.
\textbf{나} --- 무리 안 판이 판의 $5.5$ 배다.
\textbf{다} --- 짝지으면 바닥보다 $3\!\sim\!17$ 배 낫다.}
\end{figure}

\section{무리 안 판이 제일 약하다}

무리 안 판은 같은 $3{,}369$ 행을 쓰는데 표본 SD 가 \textbf{판의 $5.5$
배}다($0.0114$ 대 $0.0021$). 무리(\texttt{grp})마다 $8$ 행 이상인 칸만 쓰고
그 안에서 $\rho$ 를 내 평균하므로 \textbf{유효 표본이 훨씬 작다} --- 미리
적은 대로다.

\begin{quote}
\textbf{$t{=}0.2$ 는 이 자가 이 채택에 대해 아무 말도 못 한다는 뜻이다.}
노트 $586$ 은 무리 안 판을 ``$+0.0013$($24/40$), $0$ 을 문다''로 적고
넘어갔는데, 그것은 \emph{씨앗} 기준이었고 표본으로는 $[-0.023, +0.033]$
이다. \textbf{``$0$ 을 문다''가 같은 말이어도 폭이 $12$ 배 다르다.}
\end{quote}

\section{셋째 짝은 관문이 될 수 없다}

인계 카드가 두 달 동안 ``\textbf{CN 만화를 다섯 자에 넣을까 --- 셋째 집 밖
짝이 서면 거울(노트 $426$)이 단단해진다}''를 걸어 두고 있었다.

\begin{center}
CN 차 $+0.0105$ \quad 표본 SD $\mathbf{0.0109}$ \quad $95\%$
$[-0.0191, +0.0369]$
\end{center}

\textbf{차가 자기 잡음과 같은 크기다.} $352$ 행은 이 판의 채택 크기를
가르기에 부족하다 --- \textbf{관문으로는 안 넣는다.} 진단용으로는 남긴다
(노트 $577$ 의 천장 대비 $0.88$ 은 여전히 쓸 수 있다 --- 그것은 \emph{차}가
아니라 \emph{수준}이다).

\begin{quote}
\textbf{수준은 작은 표본에서도 말하지만 차는 못 한다.} 같은 $352$ 행이
``CN 이 전이가 되나''($+0.32$)에는 충분하고 ``이 변경이 CN 을 올렸나''
($+0.0105 \pm 0.0109$)에는 부족하다. \textbf{짝을 관문에 넣을지는 수준이
아니라 차로 정해야 한다.}
\end{quote}

\section{짝지으면 노트 $562$ 의 바닥보다 훨씬 낫다}

노트 $562$ 가 ``고르는 단위마다 검증 표본이 효과보다 커야 한다 --- 최소
크기 $\approx (2/e)^2$''를 세웠다. 그것은 \textbf{짝 안 지은} 추정의
바닥이다. 실측은

\begin{center}
\begin{tabular}{lrrr}
\hline
자 & $2/\sqrt{n}$ & 실측 SD & 배 \\
\hline
판 & $0.034$ & $0.0021$ & $\mathbf{17}$ \\
날짜 통제 판 & $0.034$ & $0.0026$ & $13$ \\
무리 안 판 & $0.034$ & $0.0114$ & $3$ \\
KR 만화 & $0.048$ & $0.0123$ & $4$ \\
비게임 앱 & $0.050$ & $0.0122$ & $4$ \\
CN 만화 & $0.107$ & $0.0109$ & $\mathbf{10}$ \\
\hline
\end{tabular}
\end{center}

\textbf{두 팔이 같은 행을 보므로 표본 흔들림이 상당 부분 상쇄된다.} 노트
$562$ 의 바닥을 그대로 쓰면 KR($+0.039$)과 CN 을 둘 다 ``못 잰다''로
버렸을 것이다 --- \textbf{KR 은 실제로는 $t{=}3.0$ 으로 잰다.}

\begin{quote}
노트 $610$ 이 \emph{씨앗} 쪽에서 본 것과 같은 구조다 --- \textbf{짝지어짐이
잡음을 깎는다.} 씨앗 쪽에서는 그것이 함정이었고(좁은 구간을 확실함으로
읽었다) 표본 쪽에서는 \textbf{이득}이다. 차이는 \emph{무엇을 다시
뽑느냐}가 아니라 \textbf{그 구간이 답하는 물음이 무엇이냐}다.
\end{quote}

\section*{조항}

\begin{quote}
\textbf{자를 관문에 넣기 전에 그 자의 차 잡음을 잰다.} 수준이 좋아도 차를
못 가르면 관문이 아니다. 그리고 \textbf{잡음은 자마다 $5.5$ 배까지
다르므로 하나의 문턱을 여섯 자에 같이 쓸 수 없다} --- 판은 $0.0045$,
무리 안 판은 $0.023$, CN 은 $0.022$ 가 각자의 $2$ SD 다.
\end{quote}

\section*{남은 것}

\begin{itemize}
\item \textbf{차 SD 는 팔 짝마다 다를 수 있다.} 노트 $610$ 이 판에서 두 팔
($0.00227 \cdot 0.00229$)로 안정함을 봤지만, \textbf{집 밖 짝에서는
안 봤다.} CN 의 $0.0109$ 가 이 팔 짝에서만의 값일 수 있다.
\item \textbf{무리 안 판을 고칠 수 있나} --- 무리 문턱 $8$ 을 올리면
잡음이 줄고 표본이 더 줄어든다. \textbf{어느 쪽이 이기는지 안 쟀다.}
\item \textbf{여섯 자의 상관을 안 쟀다} --- 판과 날짜 통제 판은 같은 행을
쓰므로 거의 같은 말을 할 것이다. \textbf{독립인 자가 몇 개인지 모른다.}
\end{itemize}

\end{document}
